\section{Introduction}
\label{sec:intro}

Progress in the security of large language model (LLM) agents has tracked the quality
of the instruments used to measure it. Agent Security Bench (ASB)~\cite{zhang2025asb}
formalized attacks and defenses across an agent's operational stages and introduced a
utility--security metric $\mathrm{NRP}=\mathrm{PNA}\cdot(1-\mathrm{ASR})$;
AgentDojo~\cite{debenedetti2024agentdojo} evaluates an agent under prompt injection
with \emph{deterministic}, state-based \texttt{security()} checks, so that a
successful injection cannot also corrupt the evaluator; InjecAgent~\cite{zhan2024injecagent}
benchmarks indirect prompt injection in tool-integrated agents; and the ToolEmu
sandbox~\cite{ruan2024toolemu} lets a language model emulate arbitrary tool execution
from specifications, trading fidelity for breadth. These instruments are strong, and
they are all single-agent. Their only cross-agent channel is a shared memory store,
and none of them exposes a knob for the phenomena that arise when several agents act
together.

Deployment has moved in the opposite direction. Production systems such as
AutoGen~\cite{wu2023autogen}, MetaGPT~\cite{hong2023metagpt},
CAMEL~\cite{li2023camel}, ChatDev~\cite{qian2024chatdev},
AgentVerse~\cite{chen2024agentverse}, and Magentic-One~\cite{fourney2024magenticone}
are orchestrated \emph{collectives}: a planner directs workers, agents read and write a
shared memory, third-party tools are invoked on the critical path, and many tasks are
structured so that no single agent can finish them alone. In such systems the security
questions that matter are group-level. Do a fraction of agents \emph{collude}? Is the
\emph{orchestrator} itself trusted, or has it been compromised? Can a
\emph{malicious tool}, a \emph{Sybil} identity, or a \emph{poisoned} entry in shared
memory spread through the collective? These are not attributes of any one agent, and
re-running a single-agent benchmark $N$ times does not create them: it produces $N$
independent agents, not a coordinated adversary, an orchestrator with a trust setting,
or a compromise that propagates over a communication topology. The recent
systematization of multi-agent security (MASEC)~\cite{schroederdewitt2025masec} makes
the point sharply---security in multi-agent systems is \emph{non-compositional}, in
that individually safe agents can compose into unsafe systems---and states plainly, in
its Section~4.7, that no multi-agent security benchmark exists, enumerating what one
must provide: heterogeneous roles, partial observability, forced coordination, a
co-evolving red/blue methodology, and group-level emergence metrics. A parallel survey
of multi-agent risks from advanced AI~\cite{hammond2025multiagent} organizes the same
territory into miscoordination, conflict, and collusion, again distinct from
single-agent mitigations. The field therefore cannot currently answer a basic
question: \emph{do single-agent defenses survive multi-agent deployment?}

\paragraph{Contribution.}
We answer the instrument gap directly. \defname{} is, to the best of our knowledge and
within the verified literature, the first environment purpose-built for the security
of multi-agent LLM systems. It is not a new attack and not a new defense; it is an
environment and an evaluation methodology whose scientific object is the controlled
variation of coordination, collusion, and orchestrator variables and the system-level
metrics that quantify their effect. \Cref{fig:architecture} previews the design: a
ToolEmu-style emulation layer supplies breadth of tools; an AgentDojo-style
deterministic checker scores system-level security predicates over the global state
trace, never placing an LLM judge on the critical path; and an orchestration layer
instantiates topology, roles, partial observability, and a parameterized adversary. We
make four concrete contributions and one empirical claim.

\begin{enumerate}[leftmargin=1.5em]
\item \textbf{An environment and task suite for multi-agent LLM security}
(\Cref{sec:system,sec:method}). Heterogeneous roles (orchestrator, workers,
tool-executors, shared-memory store), partial observability, and forced-coordination
tasks, instantiated over star, chain, tree, and mesh topologies at $3$--$50$ agents.
Forced coordination, collusion, compromise propagation, and Sybil pressure are
group-level \emph{by construction} and have no single-agent analogue.

\item \textbf{A natively parameterized adversary model} (\Cref{sec:threat}). A single
configuration vector controls the adversary fraction over agents, independent versus
colluding behaviour, static versus adaptive strategy, orchestrator compromise on/off,
malicious tools, Sybil identities, and poisoned shared memory, with a system-level
compromise objective (unsafe action, exfiltration, or coordination sabotage).

\item \textbf{A system-level metric suite} (\Cref{sec:problem,sec:method}). A
generalization $\NRPs=\PNAs\cdot(1-\ASRs)$ of ASB's metric, together with collusion
success rate, compromise-propagation rate, coordination quality, time-to-detection,
recovery rate, and cost, all measured over the collective rather than per step. Because
$\ASRs$ is scored by a deterministic checker, a successful injection cannot hijack the
evaluator---a property we prove and plan to quantify an LLM-judge ablation against (deferred
to the release phase).

\item \textbf{The \corb{} co-evolving red/blue protocol} (\Cref{sec:algorithm}). A
reusable methodology that alternates a red policy proposing attack configurations from
the parameter space and a blue policy proposing defenses, logging a trajectory of
$(\text{attack},\text{defense})$ pairs and their system-level scores.

\item \textbf{A headline degradation result} (\Cref{sec:experiments}). We measure on a
transparent, clearly-labelled synthetic agent emulator that defenses effective in the
single-agent regime \emph{degrade} under collusion and compromised-orchestrator
conditions: a per-agent guardrail's relative ASR reduction over no-defense shrinks
$41.6\%\to31.8\%\to27.7\%$ across the three contexts (\Cref{tab:degradation}). Real-backbone
replication is the release-phase programme; whether the effect replicates across models and
topologies will confirm or reframe the result.
\end{enumerate}

Because reviewers rightly discount benchmarks whose central metric is asserted rather
than justified, and because the idea's own risk analysis flags construct validity as
the primary threat, we treat the theory as part of the instrument rather than as
decoration. \Cref{sec:theory} states, under explicit assumptions and with honest
proof-status labels, five results: an axiomatic characterization that singles out the
product aggregator among all utility--security combiners and reduces it to ASB's metric
when $N=1$; an evaluator-integrity theorem showing that the deterministic checker's
verdict is invariant to adversarial message content and is therefore not hijackable,
in contrast to an LLM judge; topology-dependent bounds on compromise propagation across
the four topologies together with a strict collusion-versus-independent separation in
the environment's cascade model; a forced-coordination lower bound; and finite-sample
concentration guarantees for the estimators that fix the number of episodes the sweeps
require. None of these results asserts that any real defense is secure or that any
defense provably degrades; they characterize the \emph{measurement}, which is what a
benchmark can honestly own.

\paragraph{Scope and stance.}
\defname{} is an instrument, not a verdict. A lower system-level attack success rate or
a higher system-level NRP is a \emph{measurement under a stated configuration}, not a
certificate of security; we never equate the two. Emulation fidelity is the central
threat to validity---ToolEmu's reported human agreement is $\kappaem\approx 0.48$---and
we surface it as a first-class caveat, route scoring through the deterministic checker,
and report an emulated-versus-real-tools ablation rather than treating emulator verdicts
as ground truth. Following the trusted computing base of the source design, we take the
deterministic state checker and the environment harness to be trusted and give the
emulated tools no real side effects; real-world tool side effects, physical actuation,
and model-weight attacks are out of scope. The remainder of the paper is organized as
follows. \Cref{sec:related} situates \defname{} against single-agent benchmarks,
multi-agent frameworks and attacks, and reinforcement-learning environments;
\Cref{sec:prelim} fixes notation and recalls the ASB, AgentDojo, and ToolEmu
primitives; \Cref{sec:system,sec:threat,sec:problem} give the system model, threat
model, and problem formulation; \Cref{sec:method,sec:algorithm} present the
architecture, metric suite, and \corb{} protocol; \Cref{sec:theory} develops the
theory; \Cref{sec:experiments} specifies the evaluation protocol and the synthetic-model
results; and \Cref{sec:discussion,sec:limitations,sec:ethics,sec:conclusion} discuss
implications, limitations, ethics, and conclusions.
